% R02 calibration gold; live calibration deferred under D45.
% EXPECTED[B1][card-fidelity][4]: correct control; no finding.
% EXPECTED[B2][termination][4]: correct control; no finding.
% EXPECTED[B3][split-rigour][4]: correct control; no finding.
\documentclass[11pt]{article}
\input{preamble}
\title{R02 --- Datatypes and recursion \\ \large Written answers}
\author{dafny-verus-learn-0}
\date{}
\begin{document}
\maketitle
\section*{Problem B1}\label{ans:B1}
\begin{answer}
\textbf{Task shape:} These are 161 translated HumanEval problems. The two deliverables are a faithful specification and verified executable code, each with a proof obligation.

\textbf{Oracle:} A trusted target specification judges the implementation. The generated specification separately needs equivalence to that target across possible implementations.

\textbf{Trust:} Lean checks formal statements. The human translation from the docstring to the reference statement remains an additional assumption about what should be solved.

\textbf{Cheat surface:} Copying a computable spec or satisfying a vacuous property can create misleading success. Non-computability and reference equivalence address different versions of that problem.

\textbf{Metric:} Compilation, specification proof, implementation proof and full success are distinct columns. The given systems achieve at most a single complete task among the 161 within the reported time allowance.

\textbf{Budget:} A 600-second limit is described; reference curation also takes author and reviewer time. A universal token allocation or call count is absent from this pinned summary.

\textbf{Variance:} No confidence intervals are reported. Results near zero, public tasks and a parsing-confounded row make strong comparative claims unwarranted.

\textbf{Artifact:} CLEVER is the Lean task collection with hidden reference specs described by the card. This excerpt identifies no immutable release hash or licence to put in a reproduction manifest.
\end{answer}
\section*{Problem B2}\label{ans:B2}
\begin{answer}
Compare the regex before the input size. A recursive call on a constructor's
child already lowers the first coordinate. Star also makes a call retaining
the same regex; there the suffix must be shorter. A strictly positive split
index supplies that inequality, since suffix length is original length minus
the index. Allowing a zero index removes this reason, while removing string
length from the tuple removes the coordinate that could justify this call.
The two recursive calls therefore have distinct well-founded arguments.
\end{answer}

\section*{Problem B3}\label{ans:B3}
\begin{answer}
The accepted one-character input uses an empty prefix and its full suffix:
index zero, choosing the Eps alternative. The two-character accepted input
splits after its first character, choosing the character alternative on the
left. These boundaries are unique because the final factor accepts exactly
one character. On the rejected input, index zero leaves too much suffix,
index two leaves too little, and index one gives a left prefix accepted by
neither alternative. Those exhaust the legal indices. Each match concerns
the complete slice, so merely locating the final character does not suffice.
\end{answer}

\end{document}
